\section{Experimental setup}
\label{sec:setup}

\subsection{Datasets}
\label{sec:setup_data}

% --- from chapter3_dataset.tex 3.4.1 ---
\paragraph{Synthetic data}
Because the grouping module never sees the image, its input being a set of keypoint coordinates, joint types
and visibility flags, a rendered scene is as valid a training example as a photograph, provided the geometry of the
keypoints is realistic.
Synthetic scenes also supply labels that annotation cannot.
Every keypoint's position, person identity and visibility follow from the 3D scene,
and a joint hidden behind another person is recorded as present but occluded rather than missing.
The scenes are generated in Blender~\citep{blender}: a configurable number of mannequin meshes is placed, and the
camera, viewing angle, lighting and placement are varied procedurally.
Each scene provides the seventeen COCO keypoints per person with an image position
and a visibility state determined by ray casting against the scene geometry.
The dataset contains 400 training, 100 validation and 100 test scenes.
The number of people per scene is drawn uniformly from $K \in \{1,\dots,5\}$, which gives the grouping network
and the K-head a balanced range of crowding;
the upper limit of five covers the bulk of the target distribution, since
87.5\,\% of the COCO \texttt{val2017} images with keypoint annotations contain five people or fewer.
Depth is available from the renderer but is not used.
With exact depth the standard GATv2 baseline groups synthetic scenes almost perfectly,
0.9990 accuracy on the validation split, yet transfers to COCO worse than without it, 0.8366 against 0.8841,
because real images offer only monocular depth estimates~\citep{ranftl2020midas} of a different scale.
All reported models therefore use the position-and-visibility features of Equation~\eqref{eq:node_feature_no_depth}.

% --- from chapter3_dataset.tex 3.4.2 and chapter3_evaluation.tex 3.13.1 ---
\paragraph{COCO}
Fine-tuning uses the person-keypoint annotations of COCO \texttt{train2017}~\citep{lin2014coco},
approximately 56\,000 images after filtering.
Evaluation uses \texttt{val2017} in three views.
Ground-truth keypoints: the annotated keypoints of the 2\,307 images that form a scoreable grouping instance are
grouped directly, which isolates the grouping network from any detector.
About half of these images contain a single person, for which grouping is trivial;
the 1\,248-image subset with at least two people is therefore reported alongside as the stricter measure.
End-to-end: the annotated keypoints are replaced by detections,
on the 2\,258 images that contain at least one person with three or more visible keypoints;
an image contributes only when at least two detections can be matched to ground truth, which leaves 2\,165, 2\,196
and 2\,121 scored images for HigherHRNet, YOLO11x-pose and RTMO-L respectively.
System-level: all 5\,000 \texttt{val2017} images are used for the average-precision evaluation.

\subsection{Detectors and the joint-trained reference}
\label{sec:setup_detectors}

% --- from chapter3_integration.tex 3.12.1 ---
The primary detector is HigherHRNet-W32 at an input resolution of 512 pixels~\citep{cheng2020higherhrnet,sun2019hrnet},
with its public COCO weights, held frozen throughout.
It is chosen deliberately rather than swapped for a more recent estimator.
Its associative-embedding (AE) grouping~\citep{newell2017associative} predicts a scalar tag at every keypoint,
trained jointly with the heatmaps so that tags of the same person agree and tags of different people separate.
AE is therefore a grouping mechanism learned together with the detector on the COCO training distribution,
with access to the detector's features.
It serves as the joint-trained upper reference:
the comparison measures what is lost when grouping is separated from detection,
not how PG-GAT ranks against the state of the art in pose estimation.
Both groupings receive the same candidate keypoints, extracted by HigherHRNet's own non-maximum suppression,
and their difference is therefore attributable to the grouping stage alone.

% --- from chapter3_integration.tex 3.12.2 ---
The detector swap adds two detectors from different architectural families:
YOLO11x-pose~\citep{ultralytics2024yolo11}, a single-stage model that regresses keypoints per person instance,
and RTMO-L~\citep{lu2024rtmo}, a one-stage bottom-up estimator.
Both return keypoints already assembled into people.
Their person assignments are discarded and all keypoints above a confidence of 0.1 are pooled into one detection set,
the same threshold used for HigherHRNet, before the frozen PG-GAT and the frozen K-head are applied.
No parameter is updated and neither detector contributed training data; the swap is therefore zero-shot.
The discarded assignments are retained and scored as each detector's native grouping on the same matched keypoints.
The RTMO-L weights are the public \texttt{body7} release, trained on a multi-dataset corpus rather than on COCO alone,
which favours its native grouping.

\subsection{Metrics}
\label{sec:setup_metric}

% --- from chapter3_problem_formulation.tex 3.3.4 and chapter3_integration.tex 3.12.3 ---
\paragraph{Pose grouping accuracy}
The primary metric measures the grouping decision alone.
Cluster identifiers produced by $k$-means carry no fixed meaning;
predicted clusters are therefore first aligned with ground-truth people by the Hungarian
algorithm~\citep{kuhn1955hungarian} on their overlap.
With $\pi^\star$ the resulting one-to-one mapping,
pose grouping accuracy (PGA) is the fraction of keypoints assigned to the correct person,
\begin{equation}
    \mathrm{PGA}(\hat{\mathbf{y}}, \mathbf{y}^\star)
    = \frac{1}{N}
      \sum_{i=1}^{N}
      \mathbb{1}\!\left[
          \pi^\star(\hat{\mathbf{y}}_i) = \mathbf{y}^\star_i
      \right],
    \label{eq:pga}
\end{equation}
averaged over images.
On detected keypoints, detections are first matched to annotated keypoints of the same joint type by Hungarian
assignment with a maximum distance of 10 pixels,
and each matched detection inherits the identity of its annotation.
PGA is computed over the matched detections only, for PG-GAT and the native grouping alike,
which excludes detection recall and localisation error and compares the grouping stage on an identical denominator.
In this protocol the graph given to PG-GAT is also built from the matched detections,
and unmatched false-positive keypoints do not enter it.

% --- from chapter4_e2e.tex 4.9.4 and results.md AP@OKS section ---
\paragraph{Average precision}
Published pose estimators are ranked by average precision under the object keypoint similarity (AP@OKS),
which scores complete person instances and is affected by detection failures, localisation error and grouping alike.
A fixed change in PGA has no fixed counterpart in AP;
the two metrics answer different questions ---
PGA isolates the component that PG-GAT replaces, while AP measures the system it becomes part of.
To report the second, the HigherHRNet pipeline is also evaluated with pycocotools on all 5\,000 \texttt{val2017}
images.
In this deployable setting PG-GAT receives every pooled detection above the confidence threshold,
including false positives, and no ground truth enters inference.
Each $k$-means cluster becomes one person instance holding the highest-confidence detection of each joint type,
scored by the mean confidence over the seventeen joint slots.
The native grouping is scored on exactly the same pooled joints, and additionally on its unrestricted output.
The detector wrapper runs single-scale without flip testing, and absolute values therefore sit below
the published 0.671 for this configuration~\citep{cheng2020higherhrnet};
only comparisons within the protocol are meaningful.

\subsection{Implementation}
\label{sec:setup_impl}

% --- from chapter3_baseline.tex 3.5.3 and chapter3_evaluation.tex 3.13.2-3.13.3 ---
The networks are implemented in PyTorch Geometric~\citep{fey2019pyg};
$k$-means uses scikit-learn~\citep{scikitlearn} with ten initialisations and a fixed seed.
Synthetic training selects the checkpoint with the highest pose grouping accuracy on the synthetic validation split, evaluated every five epochs.
The three hyperparameter sweeps use Bayesian optimisation with Hyperband early termination~\citep{li2018hyperband};
the architecture and loss sweeps select on pose grouping accuracy on the synthetic validation split.
Every reported number is linked to a logged run in a Weights \& Biases project~\citep{wandb},
which records the configuration, metric history and checkpoint of the run.